\documentclass[11pt,a4paper]{article}
\usepackage[margin=21mm,headheight=15pt,headsep=7mm]{geometry}
\usepackage[T1]{fontenc}
\usepackage{lmodern,microtype}
\usepackage{enumitem}
\usepackage{needspace}
\usepackage{xurl}
\usepackage{booktabs,tabularx}
\usepackage{xcolor}
\usepackage{fancyhdr}
\usepackage{amsmath}
\usepackage[colorlinks=true,urlcolor=blue!55!black,linkcolor=blue!55!black]{hyperref}
\definecolor{ink}{HTML}{263D4E}
\definecolor{linecolor}{HTML}{BCCAD1}
\setlength{\parindent}{0pt}
\setlength{\parskip}{3.5pt}
\setlist[itemize]{leftmargin=*,topsep=2pt,itemsep=1.5pt,parsep=0pt}
\renewcommand{\arraystretch}{1.13}
\pagestyle{fancy}
\fancyhf{}
\lhead{\small Individual Project Logbook}
\rhead{\small Microplastic Mapping}
\cfoot{\small Page \thepage}
\renewcommand{\headrulewidth}{0.35pt}
\newcommand{\logentry}[3]{%
  \Needspace{13\baselineskip}%
  \vspace{8pt}%
  {\large\bfseries\color{ink}#1 \quad #2}\hfill{\small\color{ink}#3}\par
  \vspace{-3pt}{\color{linecolor}\rule{\linewidth}{.6pt}}\par\vspace{1pt}%
}
\newcommand{\result}[1]{\textbf{Outcome / decision:} #1\par}
\newcommand{\evidence}[1]{\textbf{Documentary evidence:} {\small #1}\par}
\newcommand{\nextstep}[1]{\textbf{Next action:} #1\par}
\begin{document}
\begin{center}
{\LARGE\bfseries Individual Project Logbook}\\[4pt]
{\large Mapping Microplastic Concentrations in Water Bodies}\\[4pt]
{\small Research, methodology screening and preliminary autonomous-boat design}\\[2pt]
{\small Record of discussions and produced documents: October--2 November 2026}
\end{center}
\vspace{5pt}
\begin{tabular}{@{}ll@{}}
Student name: & \rule{110mm}{.4pt}\\[6pt]
Student ID: & \rule{110mm}{.4pt}\\[6pt]
Module / group: & \rule{110mm}{.4pt}\\[6pt]
Supervisor: & \rule{110mm}{.4pt}
\end{tabular}
\vspace{9pt}

\textbf{Project brief.} Design a portable system that maps measured or estimated microplastic (MP) concentration over a river or lake surface. The system should support reasonably rapid spatial coverage, location-tagged measurements, and defensible validation. A later engineering target used in the design work was to produce a map of a $20\,\mathrm{m}\times20\,\mathrm{m}$ site within about 120 minutes.

\textbf{How I prepared this log.} I reconstructed the following dated events from this project's discussion history and its generated literature-review, screening and design files. Where an event covered several discussions, the date is shown as a range. The record distinguishes what was \emph{researched or drafted} from what remains a \emph{proposed design or experiment}. AI-assisted review and document generation are not evidence that I personally performed physical experiments. I should check the wording against my own notes and add actual hours, contributions and supervisor sign-off before submission.

\logentry{01}{Establishing the screened literature base}{Early October 2026}
\textbf{Work recorded:}
\begin{itemize}
\item I reviewed the merged Web of Science and Scopus bibliometric results: 8,923 WoS and 10,019 Scopus records, or 18,942 original records.
\item The recorded screening removed 7,318 duplicates, leaving 11,624; after excluding 727 document types, 10,897 records remained. Keyword screening classified 10,833 as preliminary include candidates and 64 for manual review.
\item I noted the four overlapping research tags: sources/causes (6,430), effects in water (7,575), laboratory/measurement methods (6,654), and data trends/mapping (5,320).
\end{itemize}
\result{The 10,833-paper pool became the starting point for more targeted sensing-method reviews. I treated these as automated screening results, not completed full-text eligibility decisions.}
\evidence{\texttt{Bibliometrix\_Filtered.RData}; original bibliometric-screening data and background LaTeX review.}
\nextstep{Narrow the literature to methods that could plausibly work without a full laboratory.}

\logentry{02}{Reviewing camera-based AI and EIS as direct sensors}{14 October 2026}
\textbf{Work recorded:}
\begin{itemize}
\item I requested and examined an in-depth comparison of camera + AI particle counting and electrical impedance spectroscopy (EIS) for a water-monitoring boat.
\item The optical route considered a pump, calibrated flow, transparent chamber, controlled lighting, imaging, image segmentation and particle counting. The EIS route considered an electrode chamber, electrical frequency response and the influence of water conductivity.
\item The resulting review used 22 DOI-matched references and compared the two principles, practical issues, sample volume and possible validation experiments.
\end{itemize}
\result{Camera imaging gave a potentially auditable particle count, while EIS offered a possible compact sensor but greater uncertainty about particle specificity in natural water. Neither was assumed to identify polymers automatically.}
\evidence{\texttt{Camera\_AI\_Impedance\_EIS\_Literature\_Review\_Overleaf.zip} and its compiled 11-page review.}
\nextstep{Explain the two methods in sufficient detail to decide what the process workflow requires.}

\logentry{03}{Clarifying what the sensors actually measure}{15 October 2026}
\textbf{Work recorded:}
\begin{itemize}
\item I compared image segmentation (including U-Net-style methods) with the complex EIS response $Z(f)=V(f)/I(f)$.
\item I separated image/particle-detection accuracy from a valid concentration measurement in particles/L. I also distinguished conventional frequency-swept EIS from DC microfluidic resistance sensing.
\item I reviewed examples including Lee et al. (2022) for image counting, Choran and \"Ormeci (2023) for micro-flow imaging, and Kim et al. (2024) for a PVC-focused electrochemical sensor.
\end{itemize}
\result{I concluded that a sensor's published image-segmentation score or electrical detection limit cannot be substituted for its accuracy at estimating MP concentration in untreated river water.}
\evidence{Discussion-based comparison and the Camera + AI versus EIS report; Lee et al., DOI: \href{https://doi.org/10.1007/s00521-021-06876-w}{10.1007/s00521-021-06876-w}.}
\nextstep{Define the overall workflows before choosing the technology.}

\logentry{04}{Defining two ways to complete the stakeholder task}{16 October 2026}
\textbf{Work recorded:}
\begin{itemize}
\item I reorganised the project into \textbf{Track A}, involving measurements on a moving boat, and \textbf{Track B}, involving geotagged water collection followed by portable shoreline processing.
\item I included common requirements for both tracks: known water volume, time and GPS records, processing quality checks, and a 2D concentration map.
\item I explicitly removed an assumption of routine FTIR/Raman laboratory access from the practical design path; occasional reference testing remained desirable if available.
\end{itemize}
\result{Both tracks could serve the stakeholder's map objective, but Track B would produce results after collecting and processing samples rather than continuously in the water.}
\evidence{\texttt{Microplastic\_Dual\_Track\_Background\_Overleaf.zip}; two-track background review.}
\nextstep{Turn the flow diagram into a criterion for selecting sensors, rather than a fixed camera-based design.}

\logentry{05}{Making the process flowchart method-neutral}{16 October 2026}
\textbf{Work recorded:}
\begin{itemize}
\item I requested that the flowchart stop specifying a camera and EIS in its process boxes.
\item Track A was rewritten as water intake/pump $\rightarrow$ flow-through or microfluidic chamber $\rightarrow$ \emph{candidate sensing method} $\rightarrow$ wastewater/flushing/storage $\rightarrow$ GPS, data logging and transmission.
\item Track B retained known-volume collection, sample identification, shoreline preparation, a separate \emph{candidate sensing method}, and linking the results back to the boat's GPS and time records.
\end{itemize}
\result{The sensing stage became the decision point for filtering literature methods by fit, portability, processing time, contamination control and mapping readiness.}
\evidence{\texttt{Microplastic\_Dual\_Track\_Background\_Overleaf\_v2.zip}; updated TikZ process chart.}
\nextstep{Create a more comprehensive screening and elimination process.}

\logentry{06}{Expanding screening beyond the original laboratory-method tags}{19 October 2026}
\textbf{Work recorded:}
\begin{itemize}
\item I requested a re-screen of the whole 10,833-paper candidate pool, not only the 6,654 papers tagged under laboratory methods.
\item The expanded keyword dictionary identified 49 methodology tags. Its audit output included 6,483 overlapping method assignments associated with 3,974 unique papers.
\item Of the 6,859 papers without these method tags, 137 were separately flagged as technical leads to inspect. I kept unmatched papers visible rather than treating them as definitively ineligible.
\item The review separated detection technologies, sampling and preparation, laboratory reference methods, indirect water-quality proxies, and map-generation methods.
\end{itemize}
\result{I used track-specific selection/elimination gates and a weighted comparison only after basic process-fit checks. A common method in published papers was not necessarily a practical boat-mounted sensor.}
\evidence{Method-selection review (Version 3); 49-tag CSV, paper-to-method index and weighted Track A/B comparison.}
\nextstep{Check whether key retained methods have credible size and concentration coverage.}

\logentry{07}{Investigating EIS particle-size coverage and electrodes}{21 October 2026}
\textbf{Work recorded:}
\begin{itemize}
\item I investigated the MP size definition of approximately $1\,\mu$m to $5\,\mathrm{mm}$ and clarified that size coverage is different from the dynamic range of concentration in particles/L.
\item I looked at Colson and Michel (2021) on a larger flow-through EIS channel, Butement et al. (2024) on smaller impedance-cytometry particles, and Aghel et al. (2024) on a three-electrode droplet device.
\item I discussed a hypothetical four-fraction arrangement: 1--10, 10--200, 200--1,000 and 1,000--5,000 $\mu$m. These were \emph{engineering bins}, not a demonstrated or minimum four-channel solution.
\item I noted that adding more electrodes does not by itself guarantee full size coverage; electrode geometry, channel dimensions, throughput and interference all matter.
\end{itemize}
\result{A single general-purpose electrode pair was not supported as a validated solution for the entire MP definition. I retained EIS as a feasibility candidate rather than a confirmed full-range instrument.}
\evidence{Colson and Michel (2021), Butement et al. (2024), and Aghel et al. (2024); primary papers reviewed in the EIS discussion.}
\nextstep{Find size-separation materials and assess whether they can be integrated without losing particles.}

\logentry{08}{Researching size filters for the EIS concept}{21--22 October 2026}
\textbf{Work recorded:}
\begin{itemize}
\item I sought available nominal cutoffs at 1,000, 200, 10 and 1 $\mu$m for the proposed particle-size fractions.
\item The shortlist included stainless-steel woven meshes for the larger cutoffs and track-etched membranes for the smaller cutoffs (for example, Merck \texttt{TCTP04700} at 10 $\mu$m and Whatman \texttt{10418712} at 1 $\mu$m).
\item I distinguished a sieve retaining particles from an EIS chamber actually counting them; retained fractions would require transfer, backflushing or separate analysis, with recovery losses measured.
\end{itemize}
\result{Coarse reusable mesh is more plausible for repeated onboard handling than fine membranes, which risk clogging in river water. The filters do not themselves solve the sensor problem.}
\evidence{Filter-catalogue comparison and EIS fractionation discussion.}
\nextstep{Check university-approved procurement routes and confirm product availability rather than assuming stock.}

\logentry{09}{Checking the university's registered-supplier routes}{21--22 October 2026}
\textbf{Work recorded:}
\begin{itemize}
\item I used the supplied registered-supplier list and compared RS Malaysia, element14, DigiKey, Cytron and the registered scientific/hardware purchasing routes.
\item RS listings included a 125 $\mu$m stainless suction strainer, a nominal 10 $\mu$m hydraulic element and a nominal 1 $\mu$m water cartridge; none was confirmed equivalent to a precision fractionation set.
\item I prioritised requesting exact-size mesh specifications from YLT Hardware/Grandtech/Multi Impact and precision membranes through LCG Scientific; supplier inventory and quotation status remained unconfirmed.
\end{itemize}
\result{I did not treat ordinary hydraulic or nominal-depth cartridges as validated analytical sieves. No purchase order or supplier confirmation is documented in the chat.}
\evidence{University registered-supplier list supplied in discussion; RS/element14 catalogue comparison.}
\nextstep{Request dimensions, pore/aperture tolerance, holder compatibility, quotations and lead times once the sensor architecture is chosen.}

\logentry{10}{Evaluating polarised-light imaging and the polygon hypothesis}{23 October 2026}
\textbf{Work recorded:}
\begin{itemize}
\item I investigated a suggested blacked-out imaging box/tube with fixed light, polarising filter(s), sample chamber and camera, followed by contour detection.
\item I checked the assumption that MPs should have fixed regular polygon outlines unlike non-plastic material. Published examples instead included plastic fibres, films, spheres and irregular fragments, while natural sand can also be angular.
\item I reviewed polarised scattering (Liu et al., 2021), image-shape measurements (Dacewicz et al., 2024), and polarisation-sensitive imaging/holography (Montandon and Nicolls, 2024; Lopera et al., 2026).
\item A proposed low-cost experiment was to compare ordinary images, parallel-polariser images and crossed-polariser images, using contour, elongation, texture and brightness changes together.
\end{itemize}
\result{I rejected \emph{polygon shape alone} as an MP identification rule. Polarised-light contrast plus multiple image features remained a testable Track B option, not polymer confirmation.}
\evidence{Liu et al., \href{https://doi.org/10.3389/fmars.2021.739683}{doi:10.3389/fmars.2021.739683}; literature discussed on 23 October.}
\nextstep{Trial known plastic and non-plastic particles under controlled light before developing an automatic classifier.}

\logentry{11}{Screening freshwater MP size studies}{30 October 2026}
\textbf{Work recorded:}
\begin{itemize}
\item I requested a size-focused screen of the 10,833 records, particularly rivers and ponds, to inform the size range that the first prototype should target.
\item The automated title/abstract screen generated 2,085 freshwater-and-size candidates; 17 studies were curated for detailed examples. I treated these 17 as selected examples, not a comprehensive full-text systematic review.
\item I compared reports from the Langat and Selangor systems and international rivers/ponds, noting examples of high frequencies of particles below 1 mm, with reported distributions depending on methods and detection cutoffs.
\end{itemize}
\result{The literature did not justify one universal ``most common'' freshwater size. A sensor's minimum measurable size must be stated and compared with the methods used in local papers.}
\evidence{Freshwater size-literature review; 2,085-paper first-pass CSV and curated size-study table.}
\nextstep{Separate what is actually known about Semenyih River from broader Malaysian distributions.}

\logentry{12}{Checking Semenyih River concentration and size evidence}{30 October 2026}
\textbf{Work recorded:}
\begin{itemize}
\item I focused the review on Semenyih River papers rather than transferring percentages from other rivers.
\item A Chen et al. (2023) study was recorded as reporting an average river-water concentration of $3.12\pm2.49$ particles/L and fibre dominance.
\item The accessible study summary did not establish reliable numerical percentages by particle-size bin. I therefore avoided producing a made-up Semenyih histogram.
\end{itemize}
\result{Semenyih-specific concentration and shape information can guide sample-volume planning, but a local size-frequency distribution remains an evidence gap until detailed figures or underlying data are checked.}
\evidence{Chen et al. (2023), \href{https://doi.org/10.1007/s10661-022-10856-5}{doi:10.1007/s10661-022-10856-5}; Semenyih-focused review.}
\nextstep{Look for a full-text size histogram or obtain more local size data before fixing the EIS/microscopy target range.}

\logentry{13}{Finding University of Nottingham Malaysia laboratory studies}{30 October 2026}
\textbf{Work recorded:}
\begin{itemize}
\item I searched affiliation metadata, separating University of Nottingham Malaysia from the Nottingham UK and Ningbo campuses. The broad Nottingham match set contained 33 records; strict Malaysia-affiliation matching produced nine.
\item I reviewed reported laboratory workflows including field-water sampling, Fenton digestion, filtration on glass microfibre media, stereomicroscopy and FTIR checks. The actual combination depended on the study and sample type.
\item One Chen et al. (2021) study reportedly checked 85 visually identified particles with FTIR, with 96\% correct visual identification in that checked subset.
\end{itemize}
\result{I kept the 96\% figure distinct from sample recovery, limits of detection and full-method concentration accuracy. A literature affiliation match does not establish that all laboratory equipment is available for this project.}
\evidence{UNM nine-paper evidence table and freshwater size/UNM literature review PDF.}
\nextstep{Use UNM's published methods as contextual reference procedures and verify equipment access separately if needed.}

\logentry{14}{Making the freshwater/UNM screening reproducible}{30 October--2 November 2026}
\textbf{Work recorded:}
\begin{itemize}
\item I requested that the literature package contain a dedicated \texttt{scripts/} folder and auditable decisions rather than only summary tables.
\item The updated files included Python screening and verification scripts, an RStudio helper to export the relevant object from \texttt{Bibliometrix\_Filtered.RData}, manual curation records and an all-record screening-flags CSV.
\item The documented rerun used the existing 10,833-record CSV export, reproducing 2,085 freshwater/size first-pass matches and nine strict Nottingham Malaysia affiliation matches. It did not claim to run R directly on the uploaded RData in that environment.
\end{itemize}
\result{The paper-selection trail was made more reproducible. Automated matching, manually curated examples and verified full-text evidence remain separate levels.}
\evidence{Reproducible freshwater/UNM screening package, Python scripts and 10,833-record flags CSV.}
\nextstep{Preserve raw records and continue full-text checking of the papers most relevant to the sensor's detectable size range.}

\logentry{15}{Adding the 20 by 20 m / two-hour survey constraint}{30 October 2026}
\textbf{Work recorded:}
\begin{itemize}
\item I added a numerical mission requirement: produce a spatial map of a 400 m$^2$ region within 120 minutes, including setup, measurement, processing and rendering.
\item Track A's model used five 20 m transects, 5 m lane spacing, 120 m nominal route distance, 0.25 m/s speed, 60 min measurement allowance and 25 min for other activities, giving a 93 min total.
\item Track B's model used a 4-by-4 grid (16 stations), 1 L/station, 1.5 L/min actual collection rate, 1.2 min handling and 3 min processing per station, giving about 95.9 min in the assumed schedule.
\item I noted the sensitivity to shoreline processing time: increasing processing from 3 to 5 min per station would exceed the two-hour target.
\end{itemize}
\result{Both tracks were only \emph{modelled} as fitting the deadline. GPS coverage, effective measured water volume and actual process time remained untested.}
\evidence{\texttt{Method\_Rationale\_20m\_2h\_Overleaf\_v4.zip}; operational feasibility section.}
\nextstep{Time sample collection and analysis with a real prototype; revise the station plan using measured throughput.}

\logentry{16}{Adding operational method-selection gates and count uncertainty}{30 October 2026}
\textbf{Work recorded:}
\begin{itemize}
\item I introduced operational gates G6--G10 for end-to-end survey time, adequate analysed volume, navigation precision, energy/flotation, and data integrity/uncertainty.
\item Using the published Semenyih mean of 3.12 particles/L as an illustrative input, 1 L would contain about 3.12 particles on average under perfect detection. The approximate Poisson relative counting error is $1/\sqrt{3.12}\approx57\%$.
\item The method review concluded that a fast boat and smooth heat map alone cannot demonstrate reliable concentration measurement; sampling volume and particle specificity may be the main bottlenecks.
\end{itemize}
\result{Track B remained the lower-risk first prototype \emph{if} rapid processing could be demonstrated. Track A retained the advantage of denser spatial sampling but required an effective-volume and counting validation.}
\evidence{\texttt{2h\_Method\_Operational\_Gates.csv}; revised Method Rationale V4.}
\nextstep{Set acceptable uncertainty and the minimum sample volume before locking in the station grid or sensing rate.}

\logentry{17}{Researching the survey boat and its subsystems}{30 October 2026}
\textbf{Work recorded:}
\begin{itemize}
\item I requested a separate critical-components LaTeX report on hull/payload, RC and telemetry links, power distribution, motor options, GPS with active antenna, controller and pump specifications, and battery runtime.
\item The proposed architecture used two pontoons and differential thrust, with a Pixhawk 6C/ArduPilot Rover boat configuration, NEO-M9N GNSS, separate manual RC override and data telemetry, and a 12 V-class sampling pump.
\item The report treated propulsion thrust as separate from buoyancy/payload; it did not use a motor's thrust rating as a direct maximum boat payload.
\end{itemize}
\result{The subsystems were organised into a provisional bill of components and design comparisons; these were research selections, not confirmed purchases or assembled hardware.}
\evidence{\texttt{Boat\_Critical\_Components\_Overleaf\_v1.zip}; \texttt{boat\_components.tex}.}
\nextstep{Choose a hull geometry and weigh all proposed hardware before purchasing propulsion and power components.}

\logentry{18}{Checking the preliminary mass, buoyancy and battery model}{30 October 2026}
\textbf{Work recorded:}
\begin{itemize}
\item The model considered two pontoons each 0.9 m long and 0.2 m in diameter; an illustrative 60\%-submergence working displacement was about 33.9 kg versus an estimated 22.0 kg loaded mass including 8 L of samples.
\item The battery example was 12.8 V, 20 Ah LiFePO$_4$: 256 Wh nominal and 204.8 Wh usable under an assumed 80\% usable-energy factor.
\item Scenario power assumptions gave roughly 4.27 h at 48 W (Track B), 2.33 h at 88 W (Track A) and 1.64 h at 125 W (more difficult water). These are calculation outputs, not measured endurance.
\item I noted that combined peak current from two T200-class thrusters could exceed an example 30 A continuous battery-management-system rating.
\end{itemize}
\result{The boat had theoretical energy and buoyancy margin in the nominal scenarios, but actual thrust/current, stability, sample-water payload and battery/BMS protection remained unverified.}
\evidence{\texttt{design\_inputs.json}; \texttt{recalculate\_boat.py}; Boat Critical Components report.}
\nextstep{Carry out weighed-payload flotation, motor current, pump-backpressure and two-hour loaded runtime tests.}

\logentry{19}{Consolidating the initial research record}{2 November 2026}
\textbf{Work recorded:}
\begin{itemize}
\item I requested a complete project logbook, then a shorter first-person individual version rather than a group research report.
\item The first logbook package recorded a 22-page retrospective project history; the individual revision condensed it into four pages with dated entries and next actions.
\item I reviewed the individual format and requested a further revision with specific events, values, papers and documentary outputs rather than broad statements such as ``researched sensors.''
\end{itemize}
\result{The logbook format is now intended to serve as an individual record of documented research and design decisions, while leaving physical build/testing entries for later.}
\evidence{22-page retrospective research logbook and four-page first-person logbook (2 November).}
\nextstep{Append separately dated, contemporaneous entries when parts are purchased, the prototype is assembled, or actual measurements are taken.}

\section*{Current status and next recorded tests}
The discussion and generated files document literature research, classification scripts, concept evaluations, a preliminary boat architecture and numerical feasibility models. They do \textbf{not} establish that a sensor was built, a laboratory measurement was performed, a component was purchased, or a $20\,\mathrm{m}\times20\,\mathrm{m}$ survey was completed.

\textbf{Next practical tasks:}
\begin{enumerate}[leftmargin=*,itemsep=2pt,topsep=2pt]
\item Select the detectable particle-size band and specify whether the first output is suspected particles/L or polymer-confirmed MPs/L.
\item Compare polarised-light imaging and/or an EIS test channel with known plastic, sediment, fibre and algae samples; log blanks, recovery and false detections.
\item Time the full Track B sampling and portable processing workflow and measure the actual observation volume for Track A.
\item Weigh the boat components, test buoyancy, motor current, pump performance and two-hour battery endurance under realistic loads.
\item Generate a map from actual georeferenced observations and test interpolation at withheld locations.
\end{enumerate}

% For a new event, copy the pattern below and replace with factual details:
% \logentry{20}{Short description}{Date}
% \textbf{Work recorded:} ...
% \result{...}
% \evidence{...}
% \nextstep{...}
\end{document}